\section{Related Work}
\label{sec:related}

\paragraph{Compositional Reasoning Benchmarks.}
A wave of benchmarks has exposed fundamental compositional weaknesses in VLMs. Winoground~\cite{thrush2022winoground} demonstrated that contrastive VLMs like CLIP struggle with visio-linguistic compositionality. ARO~\cite{yuksekgonul2023aro} showed that these models behave like ``bags of words,'' ignoring attribute-object bindings and word order. CREPE~\cite{ma2023crepe} tested systematicity and productivity of compositional reasoning, while VALSE~\cite{parcalabescu2022valse} and VL-CheckList~\cite{zhao2022vlchecklist} provided broad evaluations across linguistic phenomena. However, SugarCrepe~\cite{hsieh2023sugarcrepe} and ConMe~\cite{huang2024conme} revealed that many of these benchmarks contain exploitable biases, inflating compositionality scores. \method avoids this issue entirely: because both images in each pair are synthetically generated, there are no linguistic shortcuts to exploit.

\paragraph{Counterfactual Evaluation of VLMs.}
The idea of using counterfactual modifications for more robust evaluation has gained traction. C-VQA~\cite{zhang2024whatif} added counterfactual presuppositions to VQA questions, revealing up to 40\% performance drops. CounterCurate~\cite{zhang2024countercurate} used GPT-4V and DALL-E~3 to generate counterfactual image-text pairs for fine-tuning. Chen~\etal~\cite{chen2020counterfactual} synthesized counterfactual VQA samples to reduce language bias. CausalVLBench~\cite{komanduri2025causalvlbench} tests VLMs on structural causal inference. Our work differs from all of these in a key respect: we evaluate \emph{consistency} across paired images rather than accuracy on individual counterfactual questions. This means \method measures whether models track the causal structure of the scene, not just whether they can answer harder questions.

\paragraph{Visual Grounding and Hallucination.}
Tong~\etal~\cite{tong2024eyeswideShut} introduced MMVP, showing that visual shortcomings in CLIP encoders propagate to downstream multimodal LLMs. POPE~\cite{li2023pope} exposed systematic object hallucination in large VLMs, while Visual Contrastive Decoding~\cite{leng2024vcd} proposed a training-free mitigation. Zeng~\etal~\cite{zeng2024arpgrounding} found that VLMs fail at compositional visual grounding. Our occlusion and spatial reasoning categories directly test these failure modes.

\paragraph{Synthetic Diagnostic Benchmarks.}
The use of synthetic data for diagnostic evaluation has a rich history. CLEVR~\cite{johnson2017clevr} pioneered the use of rendered 3D scenes for compositional visual reasoning, enabling controlled evaluation of specific capabilities. GQA~\cite{hudson2019gqa} extended this to real-world images with compositional scene graphs. EqBen~\cite{wang2023eqben} tested equivariance in VLMs by modifying single elements. \method follows this tradition but introduces a new evaluation \emph{protocol}---paired counterfactual evaluation with consistency scoring---rather than just a new dataset.

\paragraph{Causal Reasoning in AI.}
Our approach is grounded in Pearl's framework of structural causal models~\cite{pearl2009causality}, which distinguishes three levels of causal reasoning: association, intervention, and counterfactuals~\cite{pearl2018bookofwhy}. Sch{\"o}lkopf~\etal~\cite{scholkopf2021causal} argued that robust AI requires causal, not merely statistical, representations. K{\i}c{\i}man~\etal~\cite{kiciman2023causal} found that LLMs have emergent but brittle causal reasoning capabilities. \method operationalizes the interventionist level of Pearl's hierarchy for VLM evaluation: we literally intervene on the visual scene and check whether the model's output reflects the intervention.
